% =====================================================================
%  DANH MỤC TOÁN 10 — tên bộ tài liệu, chương, chuyên đề.
%  Đầu trang và tiêu đề mọi tài liệu lấy tên từ đây. Đổi tên thì chỉ sửa file này.
%  Dùng trong file .tex:  \loai{1}\chuong{3}   |   \loai{2}\de{01}   |   \loai{3}\chuyende{1}
% =====================================================================
\dmuc{loai1}{CÁC CHUYÊN ĐỀ CƠ BẢN TOÁN 10}
\dmuc{loai2}{ĐỀ KIỂM TRA TOÁN 10}
\dmuc{loai3}{CHUYÊN ĐỀ ÔN THI TOÁN 10}

% Chương SGK Toán 10 Kết nối tri thức (kiểm lại theo sách đang dạy)
\dmuc{chuong1}{Chương I. Mệnh đề và tập hợp}
\dmuc{chuong2}{Chương II. Bất phương trình và hệ bất phương trình bậc nhất hai ẩn}
\dmuc{chuong3}{Chương III. Hệ thức lượng trong tam giác}
\dmuc{chuong4}{Chương IV. Vectơ}
\dmuc{chuong5}{Chương V. Các số đặc trưng của mẫu số liệu không ghép nhóm}
\dmuc{chuong6}{Chương VI. Hàm số, đồ thị và ứng dụng}
\dmuc{chuong7}{Chương VII. Phương pháp toạ độ trong mặt phẳng}
\dmuc{chuong8}{Chương VIII. Đại số tổ hợp}
\dmuc{chuong9}{Chương IX. Tính xác suất theo định nghĩa cổ điển}
% Tên ngắn cho đầu trang (chỉ cần với tên dài)
\dmucngan{chuong2}{Chương II. Bất phương trình bậc nhất hai ẩn}
\dmucngan{chuong5}{Chương V. Số đặc trưng của mẫu số liệu}

% Chuyên đề (đặt tạm, sửa theo ý)
\dmuc{chuyende1}{Chuyên đề 1. Mệnh đề, tập hợp và bất phương trình}
